% Book 3, Year 1, chapter 15 -- linked structures.
% Provenance (never printed): sources/book3/BRIEFS.md, chapter 15.
\chapter{Linked Structures: Lists, Stacks and Queues}\label{ch:b1:linked-structures}

An editor keeps the history of a long document as a sequence of a million edits, and the user goes back to the middle
of it to insert a forgotten change. In a \omterm{def:b1:dynamic-arrays:dynamic-array}{dynamic array}, inserting in the middle moves half a million entries one cell
to the right. In a \omterm{def:b1:linked-structures:list}{linked list}, it changes two \omterm{def:b1:pointers-and-arrays:pointer}{pointers}. Linked structures give up the \omterm{def:b1:pointers-and-arrays:array}{array}'s direct access to the
$i$-th element in exchange for insertions and removals that touch only the neighbours. This chapter builds them in C,
puts \omterm{def:b1:linked-structures:stack}{stacks} and \omterm{def:b1:linked-structures:stack}{queues} on top of both kinds of structure behind one interface each, separates the loop over a
structure from the structure with \omterm{def:b1:linked-structures:iterator}{iterators}, and measures what the table of costs does not say.

\begin{recall}
An \omterm{def:b1:dynamic-arrays:adt}{abstract data type} is a set of operations with a specification; its implementations are interchangeable behind its
interface (\cref{ch:b1:dynamic-arrays}). A \omterm{def:b1:pointers-and-arrays:pointer}{pointer} holds an address; \code{p->f} is the field \code{f} of the
structure that \code{p} points to; a cell obtained from \code{malloc} must be returned with \code{free} exactly once, and
never used afterwards (\cref{ch:b1:pointers-and-arrays,ch:b1:dynamic-memory}).
\end{recall}

\section{Singly linked lists}\label{sec:b1:linked-structures:singly}

\begin{definition}[Linked list, cell]\label{def:b1:linked-structures:list}
A \emph{linked list}\index{linked list} stores a sequence in \emph{cells}\index{cell (of a linked list)}, each holding
one element and a \omterm{def:b1:pointers-and-arrays:pointer}{pointer} to the next cell; the last cell's \omterm{def:b1:pointers-and-arrays:pointer}{pointer} is null, and the list is given by a \omterm{def:b1:pointers-and-arrays:pointer}{pointer} to its
first cell (null for the empty list). In a \emph{singly linked list}\index{singly linked list} each cell knows only its
successor; in a \emph{doubly linked list}\index{doubly linked list} it also knows its predecessor.
\end{definition}

\omcode{code/bachelor-1/15-linked-structures/listings/list.h}{6}{9}{The cell}\label{lst:b1:linked-structures:node}

Pushing at the front makes a new cell that points to the old first cell, and makes it the first: $\Theta(1)$, whatever
the length. Inserting after a known cell is just as cheap: the new cell takes the old successor, and the cell points to
the new one. The order of the two assignments matters: done the other way round, the rest of the list would be lost.

\begin{omfigure}
\begin{tikzpicture}[c/.style={draw, minimum width=7mm, minimum height=6mm, font=\small}, p/.style={draw, minimum
  width=4mm, minimum height=6mm}, >={Stealth[length=2mm]}]
\node[font=\small\ttfamily] (h) at (-1.3,0) {head};
\foreach \v/\x in {1/0, 2/1.9, 3/3.8} {\node[c] (v\v) at (\x,0) {\v}; \node[p, right=0pt of v\v] (n\v) {};}
\node[font=\small] (nul) at (5.7,0) {NULL};
\draw[->] (h) -- (v1); \draw[->] (n1.center) -- (v2); \draw[->] (n2.center) -- (v3); \draw[->] (n3.center) -- (nul);
\begin{scope}[yshift=-2cm]
\node[font=\small\ttfamily] (h') at (-1.3,0) {head};
\foreach \v/\x in {1/0, 2/1.9, 3/4.8} {\node[c] (w\v) at (\x,0) {\v}; \node[p, right=0pt of w\v] (m\v) {};}
\node[c, fill=omDef!20] (w9) at (3.35,-1.1) {25}; \node[p, right=0pt of w9] (m9) {};
\node[font=\small] (nul') at (6.7,0) {NULL};
\draw[->] (h') -- (w1); \draw[->] (m1.center) -- (w2); \draw[->] (m3.center) -- (nul');
\draw[->, thick, omThm] (m9.center) -- node[right, font=\scriptsize, pos=0.4] {1: new $\to$ next} (w3);
\draw[->, thick, omThm] (m2.center) -- node[left, font=\scriptsize, pos=0.6] {2: cell $\to$ new} (w9);
\draw[->, dashed, black!40] (m2.center) .. controls (3.6,0.6) .. (w3);
\end{scope}
\end{tikzpicture}
\par\smallskip{\small A \omterm{def:b1:linked-structures:list}{singly linked list} $1, 2, 3$, and the insertion of 25 after the cell holding 2: first the new
cell points to the successor, then the cell points to the new one; the dashed link is the one replaced.}
\end{omfigure}

Removal is the delicate operation, because the cell before the removed one must be changed, and the first cell has no
cell before it: the \omterm{def:b1:pointers-and-arrays:pointer}{pointer} \code{head} itself must change. The usual code has two cases. A \omterm{def:b1:pointers-and-arrays:pointer}{pointer} to a \omterm{def:b1:pointers-and-arrays:pointer}{pointer} removes
the distinction: \code{link} points to the link that leads to the cell under examination, either \code{head} or the
\code{next} field of a cell, and redirecting \code{*link} works the same way in both cases.

\omcode{code/bachelor-1/15-linked-structures/listings/list.c}{20}{28}{Removal through a \omterm{def:b1:pointers-and-arrays:pointer}{pointer} to a
link}\label{lst:b1:linked-structures:remove}

\begin{omfigure}
\begin{tikzpicture}[c/.style={draw, minimum width=7mm, minimum height=6mm, font=\small}, p/.style={draw, minimum
  width=4mm, minimum height=6mm}, >={Stealth[length=2mm]}]
\node[draw, font=\small\ttfamily, minimum height=6mm] (h) at (-1.4,0) {head};
\foreach \v/\x in {1/0, 2/1.9, 3/3.8} {\node[c] (v\v) at (\x,0) {\v}; \node[p, right=0pt of v\v] (n\v) {};}
\node[font=\small] (nul) at (5.7,0) {NULL};
\draw[->] (h) -- (v1); \draw[->] (n2.center) -- (v3); \draw[->] (n3.center) -- (nul);
\draw[->, dashed, black!40] (n1.center) -- (v2);
\node[font=\small\ttfamily, omThm] (link) at (2.3,1.3) {link};
\draw[->, thick, omThm] (link) -- (n1.north);
\draw[->, thick, omThm] (n1.center) .. controls (2.5,-1.0) and (3.2,-0.9) .. (v3.south);
\node[font=\scriptsize, anchor=west, align=left] at (6.6,0.2) {\code{*link = dead->next;}\\the field that pointed\\to 2 now points to 3};
\node[font=\scriptsize, anchor=north, align=center] at (-1.4,-0.5) {(\code{link} may also\\point here)};
\end{tikzpicture}
\par\smallskip{\small Removing 2 with a \omterm{def:b1:pointers-and-arrays:pointer}{pointer} to a link: \code{link} holds the address of the field that leads to the
cell, here the \code{next} field of the cell 1, and the field is redirected past the cell. Removing 1 would redirect
\code{head} the same way.}
\end{omfigure}

The version with cases needs the cell before the one removed, and therefore a separate branch for the first cell;
both versions give the same lists.

\omcode{code/bachelor-1/15-linked-structures/listings/remove_pp.c}{6}{21}{Removal with a special case for the first
cell}\label{lst:b1:linked-structures:remove-cases}

\omoutput[2-3]{code/bachelor-1/15-linked-structures/outputs/remove_pp.txt}{The list $1, 2, 3, 4$ after removing 1, 3
and the absent 5, by each version.}\label{out:b1:linked-structures:remove}

\omoutput{code/bachelor-1/15-linked-structures/outputs/list_demo.txt}{Pushes, an insertion after the second cell, and
removals, the first of the first cell.}\label{out:b1:linked-structures:demo}

Accessing the $i$-th element, finding a value, and appending at the end without a \omterm{def:b1:pointers-and-arrays:pointer}{pointer} to the last cell all walk
the list: $\Theta(n)$. Freeing a list must read \code{next} before freeing the cell that holds it
(\cref{exo:b1:linked-structures:5}).

\section{Doubly linked lists and sentinels}\label{sec:b1:linked-structures:sentinel}

A cell that knows its predecessor can be removed in $\Theta(1)$ when it is known, with no walk to find the cell before.
The four special cases of a \omterm{def:b1:linked-structures:list}{doubly linked list} --- empty list, first cell, last cell, only cell --- disappear with one
more idea.

\begin{definition}[Sentinel]\label{def:b1:linked-structures:sentinel}
A \emph{sentinel}\index{sentinel} is a cell that holds no element and is always present. In a circular \omterm{def:b1:linked-structures:list}{doubly linked
list} with a sentinel, the sentinel's successor is the first cell, its predecessor the last, and the empty list is the
sentinel pointing to itself both ways.
\end{definition}

\omcodes{code/bachelor-1/15-linked-structures/listings/dlist.c}{12-28}{Insertion and removal with no special
case}\label{lst:b1:linked-structures:dlist}

\begin{proposition}[No special case]\label{prop:b1:linked-structures:sentinel}
In a circular \omterm{def:b1:linked-structures:list}{doubly linked list} with a \omterm{def:b1:linked-structures:sentinel}{sentinel}, every cell, the \omterm{def:b1:linked-structures:sentinel}{sentinel} included, has a predecessor and a
successor, and $c\text{->next->prev} = c = c\text{->prev->next}$ for every cell $c$. Insertion after any cell and removal
of any cell other than the \omterm{def:b1:linked-structures:sentinel}{sentinel} preserve this invariant.
\end{proposition}

\begin{proof}
The empty list satisfies it: the \omterm{def:b1:linked-structures:sentinel}{sentinel} is its own neighbour both ways. Insertion of $n$ after $c$, with $d =
c\text{->next}$: afterwards $c \to n \to d$ and $d \to n \to c$ backwards, every other link unchanged, so the property
holds at $c$, $n$ and $d$, and elsewhere as before. Removal of $c$ between $p$ and $d$ links $p$ and $d$ to each other
both ways, which restores the property at both. No step uses a \omterm{def:b1:pointers-and-arrays:null}{null pointer}, because there are none.
\end{proof}

\begin{omfigure}
\begin{tikzpicture}[c/.style={draw, minimum width=8mm, minimum height=6mm, font=\small}, >={Stealth[length=1.8mm]}]
\node[c, fill=black!10] (s0) at (0,0) {S};
\draw[->] (s0.north east) .. controls (1.2,1.1) and (-1.2,1.1) .. (s0.north west);
\draw[->] (s0.south west) .. controls (-1.2,-1.1) and (1.2,-1.1) .. (s0.south east);
\node[font=\scriptsize] at (0,-1.3) {empty: the sentinel alone};
\begin{scope}[xshift=3.2cm]
\node[c, fill=black!10] (s) at (0,0) {S}; \node[c] (a) at (1.8,0) {1}; \node[c] (b) at (3.6,0) {2};
\node[c] (d) at (5.4,0) {3};
\foreach \u/\v in {s/a, a/b, b/d} {\draw[->] (\u.15) -- (\v.165); \draw[->] (\v.195) -- (\u.345);}
\draw[->] (d.north) .. controls (5.4,1.3) and (0,1.3) .. node[above, font=\scriptsize] {next of the last: S} (s.north);
\draw[->] (s.south) .. controls (0,-1.3) and (5.4,-1.3) .. node[below, font=\scriptsize] {prev of S: the last} (d.south);
\end{scope}
\end{tikzpicture}
\par\smallskip{\small A circular \omterm{def:b1:linked-structures:list}{doubly linked list} with a \omterm{def:b1:linked-structures:sentinel}{sentinel} S, empty and holding 1, 2, 3. Every cell has two
neighbours, so insertion and removal never test for an end.}
\end{omfigure}

\omoutput{code/bachelor-1/15-linked-structures/outputs/dlist_demo.txt}{Pushes at both ends, a removal of a known cell,
and pops at both ends.}\label{out:b1:linked-structures:dlist}

\section{Stacks, queues and deques}\label{sec:b1:linked-structures:stacks}

\begin{definition}[Stack, queue, deque]\label{def:b1:linked-structures:stack}
A \emph{stack}\index{stack (abstract type)} is an abstract type with operations push, pop, top and empty, where pop
removes the element pushed last (last in, first out). A \emph{queue}\index{queue} has put and take, where take removes the
element put first (first in, first out). A \emph{deque}\index{deque}, or double-ended queue, allows adding and removing
at both ends.
\end{definition}

Each has an \omterm{def:b1:pointers-and-arrays:array}{array} implementation and a linked one. A \omterm{def:b1:linked-structures:stack}{stack} on a \omterm{def:b1:dynamic-arrays:dynamic-array}{dynamic array} pushes and pops at the end, in amortised
$\Theta(1)$; on a \omterm{def:b1:linked-structures:list}{singly linked list}, at the front, in $\Theta(1)$. A \omterm{def:b1:linked-structures:stack}{queue} on a \omterm{def:b1:dynamic-arrays:ring}{circular buffer} was built in
\cref{ch:b1:dynamic-arrays}; on a \omterm{def:b1:linked-structures:list}{singly linked list} it keeps a \omterm{def:b1:pointers-and-arrays:pointer}{pointer} to each end, takes at the front and puts at the
rear. The \omterm{def:b1:linked-structures:list}{doubly linked list} of the previous section is a \omterm{def:b1:linked-structures:stack}{deque}.

\omcodes{code/bachelor-1/15-linked-structures/listings/queue_list.c}{10-13,23-39}{A \omterm{def:b1:linked-structures:stack}{queue} on a \omterm{def:b1:linked-structures:list}{linked list} with a \omterm{def:b1:pointers-and-arrays:pointer}{pointer}
to each end}\label{lst:b1:linked-structures:queue}

\begin{omfigure}
\begin{tikzpicture}[k/.style={draw, minimum width=6mm, minimum height=5mm, font=\scriptsize}, >={Stealth[length=1.6mm]}]
\node[font=\scriptsize\bfseries] at (1.4,1.0) {stack on an array};
\foreach \v/\i in {7/0, 2/1, 5/2} {\node[k] at (\i*0.6,0) {\v};}
\foreach \i in {3,4} {\node[k, black!40] at (\i*0.6,0) {};}
\draw[<-, omThm] (1.2,0.3) -- (1.2,0.75) node[above, font=\scriptsize] {}; \node[font=\scriptsize, omThm] at (1.2,-0.55) {top};
\begin{scope}[xshift=4.0cm]
\node[font=\scriptsize\bfseries] at (1.4,1.0) {stack on links};
\node[font=\scriptsize\ttfamily] (t) at (-0.3,0) {top};
\node[k] (a) at (0.6,0) {5}; \node[k] (b) at (1.5,0) {2}; \node[k] (c) at (2.4,0) {7};
\draw[->] (t) -- (a); \draw[->] (a) -- (b); \draw[->] (b) -- (c);
\end{scope}
\begin{scope}[xshift=8.0cm]
\node[font=\scriptsize\bfseries] at (1.4,1.0) {queue on an array (ring)};
\foreach \v/\i in {9/0, {}/1, 4/2, 6/3, 8/4} {\node[k] at (\i*0.6,0) {\v};}
\node[font=\scriptsize, omThm] at (1.2,-0.55) {front}; \node[font=\scriptsize, omCode] at (0,-0.55) {rear};
\end{scope}
\begin{scope}[xshift=12.0cm]
\node[font=\scriptsize\bfseries] at (1.2,1.0) {queue on links};
\node[k] (q1) at (0.1,0) {4}; \node[k] (q2) at (0.9,0) {6}; \node[k] (q3) at (1.7,0) {8}; \node[k] (q4) at (2.5,0) {9};
\draw[->] (q1) -- (q2); \draw[->] (q2) -- (q3); \draw[->] (q3) -- (q4);
\node[font=\scriptsize, omThm] (f) at (0.1,-0.65) {front}; \node[font=\scriptsize, omCode] (r) at (2.5,-0.65) {rear};
\draw[->] (f) -- (q1); \draw[->] (r) -- (q4);
\end{scope}
\end{tikzpicture}
\par\smallskip{\small \omterm{def:b1:linked-structures:stack}{Stacks} and \omterm{def:b1:linked-structures:stack}{queues} on \omterm{def:b1:pointers-and-arrays:array}{arrays} and on links. The \omterm{def:b1:linked-structures:stack}{stack} 7, 2, 5 (5 on top) and the \omterm{def:b1:linked-structures:stack}{queue} 4, 6, 8, 9
(4 at the front) each have two representations with the same operations and the same results.}
\end{omfigure}

The header \code{stack.h} declares an opaque \code{struct stack}; \code{stack_array.c} and \code{stack_list.c} each
define it. A bracket checker, a user of the \omterm{def:b1:linked-structures:stack}{stack}, is linked with each and gives the same answers; one \omterm{def:b1:c-modules-make-tests:unit-test}{test suite}, in a
header included by two test files, checks both implementations.

\omoutput{code/bachelor-1/15-linked-structures/outputs/stack_both.txt}{The bracket checker linked with the \omterm{def:b1:pointers-and-arrays:array}{array} \omterm{def:b1:linked-structures:stack}{stack}
and with the linked \omterm{def:b1:linked-structures:stack}{stack}.}\label{out:b1:linked-structures:both}

\section{Iterators}\label{sec:b1:linked-structures:iterators}

Summing the elements of an \omterm{def:b1:pointers-and-arrays:array}{array} and of a list are the same loop written twice: \code{for (i = 0; i < n; i++)} in one
case, \code{for (p = head; p != NULL; p = p->next)} in the other. An \omterm{def:b1:linked-structures:iterator}{iterator} separates what the loop does from how
the structure is walked.

\begin{definition}[Iterator]\label{def:b1:linked-structures:iterator}
An \emph{iterator}\index{iterator} over a structure is a value that designates a position in it, with an operation that
tells whether elements remain and an operation that returns the next element and advances. A loop written with an
iterator works on every structure that provides one.
\end{definition}

\omcode{code/bachelor-1/15-linked-structures/listings/iter.c}{7}{30}{One \omterm{def:b1:linked-structures:iterator}{iterator} type for lists and \omterm{def:b1:pointers-and-arrays:array}{arrays}, and a
sum that uses only it}\label{lst:b1:linked-structures:iter}

\omoutput{code/bachelor-1/15-linked-structures/outputs/iter.txt}{The same sum on both structures, and the callback
style.}\label{out:b1:linked-structures:iter}

The alternative is a \emph{callback}: the structure provides a function \code{list_each(head, f)} that calls \code{f}
on every element. The \omterm{def:b1:linked-structures:iterator}{iterator} lets the caller stop, resume, or walk two structures side by side; the callback keeps the
walk inside the structure's code. Either way, an \omterm{def:b1:linked-structures:iterator}{iterator} is valid only while the structure does not change: if the
cell an \omterm{def:b1:linked-structures:iterator}{iterator} points to is removed, the \omterm{def:b1:linked-structures:iterator}{iterator} holds the address of freed memory (\cref{exo:b1:linked-structures:4}).

\section{Costs, by proof and by measurement}\label{sec:b1:linked-structures:costs}

\begin{center}
\small
\begin{tabular}{lccc}
operation & \omterm{def:b1:dynamic-arrays:dynamic-array}{dynamic array} & \omterm{def:b1:linked-structures:list}{singly linked list} & \omterm{def:b1:linked-structures:list}{doubly linked list} \\ \hline
access the $i$-th element & $\Theta(1)$ & $\Theta(i)$ & $\Theta(i)$ \\
add or remove at the front & $\Theta(n)$ & $\Theta(1)$ & $\Theta(1)$ \\
add or remove at the end & amortised $\Theta(1)$ & $\Theta(1)$ to add (rear \omterm{def:b1:pointers-and-arrays:pointer}{pointer}) & $\Theta(1)$ \\
insert after a known position & $\Theta(n)$ & $\Theta(1)$ & $\Theta(1)$ \\
remove a known cell & $\Theta(n)$ & $\Theta(n)$ (find the one before) & $\Theta(1)$ \\
memory per element & 1 to 4 \omterm{def:b1:linked-structures:list}{cells} & element + 1 \omterm{def:b1:pointers-and-arrays:pointer}{pointer} & element + 2 \omterm{def:b1:pointers-and-arrays:pointer}{pointers}
\end{tabular}
\end{center}

The table counts operations. A traversal of $n$ elements is $\Theta(n)$ on every structure, and yet:

\omoutput{code/bachelor-1/15-linked-structures/outputs/traverse.txt}{Summing $2^{22}$ integers on the author's laptop:
in an \omterm{def:b1:pointers-and-arrays:array}{array}, in a list whose \omterm{def:b1:linked-structures:list}{cells} follow each other in memory, and in the same list after shuffling the order of its
\omterm{def:b1:linked-structures:list}{cells}.}\label{out:b1:linked-structures:traverse}

The three loops perform the same number of additions. The \omterm{def:b1:pointers-and-arrays:array}{array} is read sequentially, and the processor fetches memory
in blocks ahead of use; a list whose \omterm{def:b1:linked-structures:list}{cells} are in order is nearly as regular, with a \omterm{def:b1:pointers-and-arrays:pointer}{pointer} to follow at each step; the
shuffled list sends each step to an unpredictable address, and most of them miss every cache, a factor of about 200.
The memory hierarchy behind this is the subject of the second year. A linked structure pays for its flexibility with
locality, and its costs in the table are only a part of the story.

\begin{safety}
Every linked structure in C is a set of \omterm{def:b1:linked-structures:list}{cells} to free exactly once. A cell freed while another \omterm{def:b1:pointers-and-arrays:pointer}{pointer} still refers to
it, or a \code{next} field read after its cell was freed, is \omterm{def:b1:c-types-and-control:ub}{undefined behaviour} that may run correctly for years; the
address sanitizer (\cref{ch:b1:c-modules-make-tests}) catches it in the tests.
\end{safety}

\section{Exercises}\label{sec:b1:linked-structures:exercises}

\begin{exercise}[$\star$]\label{exo:b1:linked-structures:1}
Starting from the empty list, draw the list after: push 3, push 8, push 1, insert 5 after the first cell, remove 8, push
2, remove 2.
\end{exercise}

\begin{exercise}[$\star$]\label{exo:b1:linked-structures:2}
A \omterm{def:b1:linked-structures:stack}{stack} and a \omterm{def:b1:linked-structures:stack}{queue} start empty and receive the same operations: add 1, add 2, add 3, remove, add 4, remove, remove.
Which elements does each return, and what does each hold at the end?
\end{exercise}

\begin{exercise}[$\star$]\label{exo:b1:linked-structures:3}
What does the bracket checker of \cref{out:b1:linked-structures:both} report for \code{f(a[0)]}, \code{{}[]()} and
\code{]}? Give the \omterm{def:b1:linked-structures:stack}{stack}'s contents at the error.
\end{exercise}

\begin{exercise}[$\star$]\label{exo:b1:linked-structures:4}
An \omterm{def:b1:linked-structures:iterator}{iterator} over the list $1, 2, 3, 4$ has returned 1 and 2. What happens if the cell holding 3 is now removed and the
loop continues? And if the cell holding 2 is removed instead?
\end{exercise}

\begin{exercise}[$\star\star$]\label{exo:b1:linked-structures:5}
This function frees a list. Read the sanitizer's report, explain the error, and correct the loop.
\omcode{code/bachelor-1/15-linked-structures/listings/exo_free_bug.c}{10}{12}{Freeing a list}
\end{exercise}

\begin{exercise}[$\star\star$]\label{exo:b1:linked-structures:6}
This \omterm{def:b1:linked-structures:stack}{queue} loses its elements after it becomes empty once. Find the case it forgets, state the \omterm{def:b1:dynamic-arrays:rep-invariant}{representation
invariant} it breaks, and repair it.
\omcodes{code/bachelor-1/15-linked-structures/listings/exo_queue_bug.c}{12-35}{A \omterm{def:b1:linked-structures:stack}{queue}}
\omoutput[2-3]{code/bachelor-1/15-linked-structures/outputs/exo_queue_bug.txt}{Two elements put, and none at the
front.}
\end{exercise}

\begin{exercise}[$\star\star$]\label{exo:b1:linked-structures:7}
Write \code{list_reverse(head)}, which reverses a \omterm{def:b1:linked-structures:list}{singly linked list} in place in one pass and returns the new first cell.
State its \omterm{def:b1:program-correctness:invariant}{loop invariant}, and test it on lists of length 0, 1 and 6.
\end{exercise}

\begin{exercise}[$\star\star$]\label{exo:b1:linked-structures:8}
Write a \omterm{def:b1:linked-structures:stack}{deque} on a \omterm{def:b1:dynamic-arrays:ring}{circular buffer} of fixed capacity: push and pop at both ends in $\Theta(1)$. Test it against the
\omterm{def:b1:linked-structures:list}{doubly linked list}.
\end{exercise}

\begin{exercise}[$\star\star$]\label{exo:b1:linked-structures:9}
Prove that $n$ appends at the end of a \omterm{def:b1:linked-structures:list}{singly linked list} cost $\Theta(n^2)$ cell visits without a \omterm{def:b1:pointers-and-arrays:pointer}{pointer} to the last
cell, and $\Theta(n)$ with one.
\end{exercise}

\begin{exercise}[$\star\star\star$]\label{exo:b1:linked-structures:10}
A list built by mistake may run into a cycle: its last cell points back to an earlier one. Write \code{has_cycle(head)},
which uses two \omterm{def:b1:pointers-and-arrays:pointer}{pointers}, one moving one cell at a time and one moving two, and $O(1)$ memory. Test it on lists without
a cycle and on every cycle of lists up to length 30.
\end{exercise}

\begin{exercise}[$\star\star\star$]\label{exo:b1:linked-structures:11}
A list with a cycle has a tail of $\mu$ \omterm{def:b1:linked-structures:list}{cells} before the cycle and a cycle of $\lambda \geq 1$ \omterm{def:b1:linked-structures:list}{cells}. Prove that the two
\omterm{def:b1:pointers-and-arrays:pointer}{pointers} of \cref{exo:b1:linked-structures:10} meet after at most $\mu + \lambda$ steps of the slow \omterm{def:b1:pointers-and-arrays:pointer}{pointer}.
\end{exercise}

\begin{exercise}[$\star\star\star$]\label{exo:b1:linked-structures:12}
A \omterm{def:b1:linked-structures:stack}{queue} can be made of two \omterm{def:b1:linked-structures:stack}{stacks}: put pushes on \code{in}; take pops from \code{out}, first moving all of \code{in}
onto \code{out} if \code{out} is empty. Prove that the \omterm{def:b1:linked-structures:stack}{queue} is correct, and that $m$ operations cost $O(m)$ pushes and
pops in all.
\end{exercise}

\section{Problem: The Josephus circle}\label{sec:b1:linked-structures:problem}

\begin{problem}[{Weekend problem --- The Josephus circle}]\label{pb:b1:linked-structures:1}
$n$ people, numbered 1 to $n$, stand in a circle. Counting starts at person 1, and every $k$-th person leaves the circle;
counting goes on from the next person, until one remains. The problem is named after the historian Flavius Josephus,
from a story told about him. A program simulates the circle with a circular \omterm{def:b1:linked-structures:list}{linked list}, computes the survivor by a
recurrence, and compares both with a closed form for $k = 2$.
\omcode{code/bachelor-1/15-linked-structures/listings/josephus.c}{11}{35}{Simulation, recurrence and closed
form}\label{lst:b1:linked-structures:josephus}
\omoutput{code/bachelor-1/15-linked-structures/outputs/josephus.txt}{The order of departures for $n = 10$, $k = 3$; the
three methods for $k = 2$; and the checks.}\label{out:b1:linked-structures:josephus}

\medskip\noindent\textbf{Part I --- Simulation.}
\begin{enumerate}
\item Follow the simulation for $n = 5$, $k = 2$: the departures and the survivor.
\item Why does \code{simulate} keep a \omterm{def:b1:pointers-and-arrays:pointer}{pointer} to the person \emph{before} the next one to count?
\item Why can the \omterm{def:b1:linked-structures:list}{cells} live in one \omterm{def:b1:pointers-and-arrays:array}{array}, allocated once, although people leave?
\item Show that the simulation costs $\Theta(nk)$ steps.
\item What would the cost be with a \omterm{def:b1:dynamic-arrays:dynamic-array}{dynamic array}, removing each departing person by shifting the others?
\item Why is the circular list, rather than a list with an end, the natural structure here?
\end{enumerate}

\medskip\noindent\textbf{Part II --- The recurrence.}
\begin{enumerate}[resume]
\item Number the people from 0 to $n - 1$, and let $J(n, k)$ be the survivor's number. Show that the first to leave is
$(k - 1) \bmod n$.
\item After that departure, renumber the $n - 1$ remaining people from 0, starting at the person after the one who
left. Show that the old number of a person is (new number $+ k) \bmod n$.
\item Deduce $J(1, k) = 0$ and $J(n, k) = (J(n - 1, k) + k) \bmod n$.
\item Compute $J(n, 3)$ for $n \leq 10$ and check the survivor of the first line of the output (numbered from 1).
\item What does the recurrence cost, and why is it so much cheaper than the simulation for large $k$?
\item Check in the output that the recurrence and the simulation agree for $n \leq 1000$ and $k \leq 7$.
\end{enumerate}

\medskip\noindent\textbf{Part III --- The closed form for $k = 2$.}
\begin{enumerate}[resume]
\item With people numbered from 1, show that for $n = 2m$, after the first round the $m$ people left are $1, 3, \dots,
2m - 1$, and counting starts again at 1; deduce $J(2m) = 2J(m) - 1$.
\item For $n = 2m + 1$, the first round removes $2, 4, \dots, 2m$ and then 1; deduce $J(2m + 1) = 2J(m) + 1$.
\item With $J(1) = 1$, compute $J(n)$ for $n \leq 16$ and compare with the table.
\item Write $n = 2^a + \ell$ with $0 \leq \ell < 2^a$. Prove by strong induction that $J(n) = 2\ell + 1$.
\item What is $J(2^a)$? Explain it directly with the rounds.
\item In binary, show that $J(n)$ is $n$ rotated left by one bit (the leading 1 moved to the end).
\end{enumerate}

\medskip\noindent\textbf{Part IV --- Checking.}
\begin{enumerate}[resume]
\item The program checks the closed form against the simulation for $n \leq 10\,000$. Estimate the number of steps.
\item Why is a simulation a good oracle here, and what does the check add to the proof?
\item An independent test simulates the circle with an \omterm{def:b1:pointers-and-arrays:array}{array} of flags instead of a list. Why test with a second
simulation of a different kind?
\item For which $n$ is the survivor person $n$ itself, for $k = 2$?
\item For $k = 2$, how does $J(n + 1)$ compare with $J(n)$? Read it on the table.
\item State the named result of the problem, \emph{Josephus for $k = 2$}, with its proof and its check.
\end{enumerate}
\end{problem}
